% ═══════════════════════════════════════════════════════════════════════════════
%  CONCLUSION ET BIBLIOGRAPHIE
% ═══════════════════════════════════════════════════════════════════════════════
\newpage
\section*{Conclusion Générale}
\addcontentsline{toc}{section}{Conclusion Générale}
\label{sec:conclusion_generale}

Ce projet de fin d'études a permis la conception et le développement d'une solution complète pour la gestion des soutenances et des jurys au sein de la Faculté des Sciences Ben M'Sick. L'application, articulée autour d'une architecture moderne (Spring Boot / React / MySQL), répond aux problématiques de coordination, d'automatisation des flux documentaires et de gestion des contraintes liées à l'organisation des examens académiques.

Les objectifs initiaux, notamment l'automatisation de la planification, la détection intelligente des conflits et la génération des documents officiels, ont été atteints. Ce système constitue un levier d'amélioration opérationnelle pour le département MIP, permettant de réduire la charge administrative tout en augmentant la fiabilité et la traçabilité des processus.

Les perspectives d'évolution sont nombreuses : intégration avec le système d'information de l'université (CNE), développement d'une application mobile pour les enseignants, et amélioration des algorithmes de planification automatique basés sur l'intelligence artificielle pour une gestion encore plus fine des ressources.

\newpage
\begin{thebibliography}{99}
\addcontentsline{toc}{section}{Bibliographie}

\bibitem{spring}
Spring Boot Documentation,
\url{https://docs.spring.io/spring-boot/docs/current/reference/html/}

\bibitem{react}
React Documentation,
\url{https://react.dev/reference/react}

\bibitem{shadcn}
Shadcn/ui Documentation,
\url{https://ui.shadcn.com/}

\bibitem{tailwind}
Tailwind CSS Documentation,
\url{https://tailwindcss.com/docs}

\bibitem{git}
Pro Git Book,
\url{https://git-scm.com/book/}

\end{thebibliography}
